%==============================================================================
% Chapitre 36 - La formule de Greenhill
%==============================================================================

\section{Introduction}

À la fin du XIX\ieme\ siècle, les projectiles deviennent progressivement plus
allongés.

Les ingénieurs cherchent alors à déterminer le pas de rayure nécessaire pour
assurer leur stabilité.

Parmi les premières méthodes développées figure la formule de Greenhill.

Cette formule demeure aujourd'hui l'une des plus connues du domaine.

\begin{aretenir}

La formule de Greenhill permet d'estimer le pas de rayure nécessaire à la
stabilisation d'un projectile.

\end{aretenir}

\section{Sir Alfred George Greenhill}

La formule fut développée par Alfred George Greenhill
(1847--1927), mathématicien britannique, à la fin du XIX\ieme\ siècle.

Ses travaux visaient à répondre à une question simple :

\begin{quote}
Quel pas de rayure faut-il choisir pour stabiliser un projectile donné ?
\end{quote}

\section{Principe général}

La formule repose sur une observation fondamentale :

\begin{itemize}
    \item plus un projectile est long ;
    \item plus il nécessite de rotation ;
    \item donc plus le pas doit être court.
\end{itemize}

Cette idée reste valable aujourd'hui.

\section{Forme classique}

La formule de Greenhill s'écrit généralement :

\begin{equation}
T=\frac{C D^2}{L}
\label{eq:greenhill}
\end{equation}

où :

\begin{itemize}
    \item $T$ représente le pas de rayure ;
    \item $D$ représente le diamètre du projectile ;
    \item $L$ représente sa longueur ;
    \item $C$ est une constante empirique.
\end{itemize}

\section{La constante de Greenhill}

La valeur historique la plus connue est :

\begin{equation}
C = 150
\end{equation}

pour les vitesses classiques de l'époque.

Certaines variantes utilisent :

\begin{equation}
C = 180
\end{equation}

pour les projectiles plus rapides.

\section{Interprétation physique}

L'équation~\ref{eq:greenhill} montre immédiatement que :

\begin{itemize}
    \item un projectile plus long nécessite un pas plus court ;
    \item un projectile plus gros peut être stabilisé avec un pas plus long.
\end{itemize}

Ces tendances correspondent à l'expérience pratique.

\section{Influence de la longueur}

La longueur apparaît au dénominateur.

Si la longueur double :

\begin{itemize}
    \item le pas nécessaire est divisé par deux ;
    \item la rotation requise augmente fortement.
\end{itemize}

La longueur constitue donc un paramètre majeur.

\section{Influence du diamètre}

Le diamètre apparaît au carré.

Son influence est importante.

Cependant, dans la pratique, les variations de longueur jouent souvent un rôle
plus déterminant.

\section{Exemple simplifié}

Considérons un projectile :

\begin{itemize}
    \item diamètre : 0,308 pouces ;
    \item longueur : 1,2 pouces ;
    \item constante : 150.
\end{itemize}

La formule donne :

\begin{equation}
T=
\frac{150\times0,308^2}
     {1,2}
\approx11,9
\end{equation}

On obtient donc un pas voisin de :

\begin{equation}
1:12
\end{equation}

pouces.

\section{Pourquoi cette formule fonctionne-t-elle ?}

Greenhill a développé sa formule à partir :

\begin{itemize}
    \item d'observations expérimentales ;
    \item d'analyses mécaniques ;
    \item d'essais sur des projectiles réels.
\end{itemize}

Elle constitue donc une loi empirique.

\section{Limites de la formule}

La formule présente plusieurs limitations.

Elle ne prend notamment pas en compte :

\begin{itemize}
    \item la forme exacte du projectile ;
    \item les matériaux ;
    \item la répartition des masses ;
    \item les conditions atmosphériques ;
    \item les effets transsoniques.
\end{itemize}

\section{Projectiles modernes}

Les projectiles actuels sont souvent :

\begin{itemize}
    \item plus longs ;
    \item plus profilés ;
    \item plus complexes.
\end{itemize}

Les hypothèses de Greenhill deviennent alors moins précises.

\section{Pourquoi l'étudier encore ?}

Malgré ses limites, la formule reste utile.

Elle permet :

\begin{itemize}
    \item une estimation rapide ;
    \item une compréhension intuitive ;
    \item une première vérification.
\end{itemize}

\section{Comparaison avec les méthodes modernes}

Les méthodes modernes utilisent davantage de paramètres.

Elles fournissent généralement des résultats plus précis.

Cependant, leurs conclusions restent souvent cohérentes avec les tendances
prédites par Greenhill.

\section{Place dans l'histoire de la balistique}

La formule de Greenhill représente une étape importante de l'histoire de la
balistique.

Elle a influencé pendant des décennies :

\begin{itemize}
    \item la conception des canons ;
    \item le choix des projectiles ;
    \item les essais balistiques.
\end{itemize}

\section{Vers la formule de Miller}

La formule de Greenhill fournit une première approximation.

Les besoins modernes ont conduit au développement de méthodes plus précises.

L'une des plus connues est la formule de Miller.

Elle introduit davantage de paramètres et permet une meilleure estimation de
la stabilité gyroscopique.

\section{À retenir}

\begin{aretenir}

\begin{itemize}
    \item La formule de Greenhill estime le pas de rayure nécessaire.
    \item La longueur du projectile est un paramètre majeur.
    \item La formule est empirique.
    \item Elle reste utile pour des estimations rapides.
    \item Les méthodes modernes offrent une précision supérieure.
\end{itemize}

\end{aretenir}
